\documentclass[letterpaper,journal]{IEEEtran}
% Archived fixed Emotion-Wheel draft before the RGRPL-Base revision.

\usepackage{amsmath}
\usepackage{amsfonts}

\hyphenation{op-tical net-works semi-conduc-tor}

\begin{document}

\title{Poincar\'e Emotion-Wheel Regularization for Multimodal Emotion Recognition in Conversation}
\author{}
\maketitle

\begin{abstract}
\end{abstract}

\section{Introduction}
\label{sec:introduction}

\section{Related Work}
\label{sec:related_work}

\section{Problem Definition and Proposed Method}
\label{sec:problem_and_method}

\subsection{Problem Definition}
\label{subsec:problem_definition}

\subsubsection{Multimodal Emotion Recognition in Conversation}

Let a conversational dataset be denoted by
\(\mathcal{D}=\{\mathcal{C}_{n}\}_{n=1}^{N}\), where the \(n\)-th
conversation contains an ordered sequence of utterances
\(\mathcal{C}_{n}=\{u_{n,t}\}_{t=1}^{L_n}\).
Each utterance is associated with aligned textual, acoustic, and visual
observations,
\begin{equation}
    x_{n,t}=\left(x_{n,t}^{T},x_{n,t}^{A},x_{n,t}^{V}\right),
    \qquad
    y_{n,t}\in\mathcal{Y}=\{1,\ldots,C\},
    \label{eq:problem_input}
\end{equation}
where \(y_{n,t}\) is the ground-truth emotion class and
\(C=|\mathcal{Y}|\). The modality set is denoted by
\(\mathcal{M}=\{T,A,V\}\). The objective is to learn a contextual classifier
\(f_{\Theta}\) that predicts the emotion of every valid utterance while
exploiting both conversational context and complementary cross-modal evidence:
\begin{equation}
    \hat{y}_{n,t}
    =\underset{c\in\mathcal{Y}}{\arg\max}\;
    p_{\Theta}\!\left(c\mid
    x_{n,1:L_n}^{T},x_{n,1:L_n}^{A},x_{n,1:L_n}^{V}\right).
    \label{eq:problem_prediction}
\end{equation}
Padding utterances introduced by mini-batching are excluded from all losses
and pairwise statistics through a binary validity mask.

The standard categorical objective treats the emotion classes as independent
symbols. This assumption is restrictive for emotion recognition: confusing
two semantically adjacent emotions should not be considered geometrically
equivalent to confusing two emotions located at opposite ends of an affective
continuum. We therefore seek a model that satisfies two requirements
simultaneously. First, it must retain the discriminative and contextual
modeling ability of the original SDT architecture. Second, its modality-level
representations should respect a known circular organization of emotion
classes. The latter requirement is introduced as an auxiliary geometric prior,
rather than as a replacement for the task classifier.

\subsubsection{SDT Prediction and Training Objective}

For every valid utterance, the SDT backbone produces a pure self-modal
representation and an interaction-enhanced representation for each modality:
\begin{equation}
    h_{n,t}^{m}\in\mathbb{R}^{d_h},
    \qquad
    \widetilde{h}_{n,t}^{m}\in\mathbb{R}^{d_h},
    \qquad m\in\mathcal{M}.
    \label{eq:sdt_representations}
\end{equation}
In the implementation, the pure representations correspond to
\(H_{TT},H_{AA},H_{VV}\), whereas
\(\widetilde{h}^{T},\widetilde{h}^{A},\widetilde{h}^{V}\) correspond to
\(H'_{T},H'_{A},H'_{V}\). The enhanced features are combined by the learned
SDT fusion gate to obtain a multimodal representation
\(\widetilde{h}_{n,t}\). Specifically, the shared fusion layer produces a
feature-wise score
\(a_{n,t}^{m}=G\widetilde{h}_{n,t}^{m}\in\mathbb{R}^{d_h}\), followed by a
softmax across modalities at each feature coordinate:
\begin{align}
    [\alpha_{n,t}^{m}]_k
    &=\frac{\exp([a_{n,t}^{m}]_k)}
    {\sum_{q\in\mathcal{M}}\exp([a_{n,t}^{q}]_k)},
    \label{eq:sdt_fusion_weights}\\
    \widetilde{h}_{n,t}
    &=\sum_{m\in\mathcal{M}}
    \alpha_{n,t}^{m}\odot\widetilde{h}_{n,t}^{m},
    \label{eq:sdt_fusion}\\
    p_{n,t}
    &=\operatorname{softmax}\!\left(W_o\widetilde{h}_{n,t}+b_o\right).
    \label{eq:sdt_output}
\end{align}
Here, \(k\in\{1,\ldots,d_h\}\), and \(\odot\) denotes element-wise
multiplication. Each modality additionally has a student classifier producing
\(p_{n,t}^{m}\). The original SDT objective is
\begin{equation}
    \mathcal{L}_{\mathrm{SDT}}
    =\gamma_1\mathcal{L}_{\mathrm{task}}
    +\gamma_2\mathcal{L}_{\mathrm{student}}
    +\gamma_3\mathcal{L}_{\mathrm{KD}},
    \label{eq:sdt_objective}
\end{equation}
where \(\mathcal{L}_{\mathrm{task}}\) is the class-weighted negative
log-likelihood of the fused prediction. Defining
\(Z_{\mathcal{I}}=\sum_{i\in\mathcal{I}}\omega_{y_i}\), it is written as
\begin{equation}
    \mathcal{L}_{\mathrm{task}}
    =-\frac{1}{Z_{\mathcal{I}}}
    \sum_{i\in\mathcal{I}}\omega_{y_i}\log p_i[y_i],
    \label{eq:sdt_task_loss}
\end{equation}
where \(\mathcal{I}\) is the set of valid utterances in the mini-batch.
The unimodal supervision term is
\begin{equation}
    \mathcal{L}_{\mathrm{student}}
    =-\sum_{m\in\mathcal{M}}
    \frac{1}{Z_{\mathcal{I}}}
    \sum_{i\in\mathcal{I}}
    \omega_{y_i}\log p_i^{m}[y_i].
    \label{eq:sdt_student_loss}
\end{equation}
Finally, \(\mathcal{L}_{\mathrm{KD}}\) distills the fused prediction into the
three modality-specific students:
\begin{equation}
    \mathcal{L}_{\mathrm{KD}}
    =\sum_{m\in\mathcal{M}}
    \frac{1}{|\mathcal{I}|}
    \sum_{i\in\mathcal{I}}
    D_{\mathrm{KL}}\!\left(
    p_i^{(\tau_s)}\,\Vert\,p_i^{m,(\tau_s)}\right),
    \label{eq:sdt_kd_loss}
\end{equation}
where \(\tau_s\) is the SDT distillation temperature. We retain this prediction
path and augment only its training objective with the structured emotion prior
described next.

\subsection{Proposed Method}
\label{subsec:proposed_method}

\subsubsection{Method Overview}

We propose a fixed Poincar\'e Emotion-Wheel regularization framework for SDT.
The method attaches a geometry-aware auxiliary branch to the pure
modality-specific representations in \eqref{eq:sdt_representations}, while the
enhanced representations continue to drive the original SDT fusion and final
classifier. The auxiliary branch consists of three components: 1) a learned
projector for each modality, mapping Euclidean features into a common
\(D\)-dimensional Poincar\'e ball; 2) a shared bank of fixed class prototypes
whose angular arrangement follows the Emotion Wheel; and 3) two complementary
objectives that impose point-to-prototype discrimination and pairwise
structure preservation, respectively.

This separation is deliberate. The Emotion Wheel acts as a semantic
regularizer during optimization, but its distance-based logits are not fused
with the SDT logits. Consequently, the proposed method can shape the learned
representation without forcing an imperfect unimodal geometric classifier to
directly alter the multimodal decision at inference time.

\subsubsection{Modality-Specific Projection into the Poincar\'e Ball}

We use the unit-curvature Poincar\'e ball
\begin{equation}
    \mathbb{D}^{D}
    =\left\{z\in\mathbb{R}^{D}:\lVert z\rVert_2<1\right\},
    \label{eq:poincare_ball}
\end{equation}
whose constant sectional curvature is \(-1\). For each modality
\(m\in\mathcal{M}\), a learned affine projector first maps a pure SDT feature
to the tangent space at the origin:
\begin{equation}
    v_i^{m}=W_m h_i^{m}+b_m,
    \qquad v_i^{m}\in T_{0}\mathbb{D}^{D}\simeq\mathbb{R}^{D}.
    \label{eq:tangent_projection}
\end{equation}
The tangent vector is then mapped into the ball through the exponential map at
the origin:
\begin{equation}
    z_i^{m}=\operatorname{Exp}_{0}(v_i^{m})
    =\tanh\!\left(\lVert v_i^{m}\rVert_2\right)
    \frac{v_i^{m}}{\lVert v_i^{m}\rVert_2},
    \label{eq:poincare_expmap}
\end{equation}
with the zero-vector case defined by continuity. This parameterization
guarantees \(z_i^{m}\in\mathbb{D}^{D}\). Importantly, the radius of each sample
is learned freely; no normalization forces all utterances onto a common
hypersphere. Thus, the representation may encode information jointly in its
angular direction and radial magnitude.

For \(x,y\in\mathbb{D}^{D}\), their geodesic distance is
\begin{equation}
    d_{\mathbb{D}}(x,y)
    =\operatorname{arcosh}\!\left(
    1+2\frac{\lVert x-y\rVert_2^2}
    {(1-\lVert x\rVert_2^2)(1-\lVert y\rVert_2^2)}
    \right).
    \label{eq:poincare_distance}
\end{equation}
Unlike Euclidean distance, \eqref{eq:poincare_distance} expands rapidly near
the boundary of the ball, providing a geometry in which angular class
organization and sample-dependent radial variation can coexist.

\subsubsection{Fixed Emotion-Wheel Prototype Bank}

Let \(\pi=(\pi_0,\ldots,\pi_{C-1})\) denote the cyclic ordering of emotion
classes supplied by the Emotion Wheel. We assign uniformly spaced angles
according to
\begin{equation}
    \theta_{\pi_k}=\frac{2\pi k}{C},
    \qquad k=0,\ldots,C-1.
    \label{eq:wheel_angles}
\end{equation}
For IEMOCAP, the adopted cyclic order is
\begin{equation}
\begin{aligned}
    \textit{happy}&\rightarrow\textit{excited}\rightarrow
    \textit{angry}\rightarrow\textit{frustrated}\\
    &\rightarrow\textit{sad}\rightarrow\textit{neutral}
    \rightarrow\textit{happy}.
\end{aligned}
    \label{eq:iemocap_wheel_order}
\end{equation}
Each emotion class \(c\) is represented by a prototype at a common radius
\(r\in(0,1)\):
\begin{equation}
    q_c=
    \left[r\cos\theta_c,\;r\sin\theta_c,\;
    \underbrace{0,\ldots,0}_{D-2}\right]^{\top}
    \in\mathbb{D}^{D}.
    \label{eq:fixed_wheel_prototype}
\end{equation}
The bank \(\mathcal{Q}=\{q_c\}_{c=1}^{C}\) is shared by text, audio, and
visual modalities. Its coordinates are fixed throughout training; only the SDT
encoder and the modality projectors are optimized. The shared bank therefore
provides a stationary semantic reference frame and prevents the three
modality-specific projectors from defining mutually inconsistent class
geometries.

Although the prototypes occupy a two-dimensional semantic plane, the sample
representations \(z_i^{m}\) remain fully \(D\)-dimensional. In particular, the
coordinates from 3 to \(D\) are not truncated. They contribute to both
\eqref{eq:poincare_distance} and the pairwise loss introduced below, and hence
receive gradients during training. The construction should therefore be
interpreted as a \(D\)-dimensional feature space regularized by a shared
two-dimensional semantic plane, rather than as a purely two-dimensional
embedding.

\subsubsection{Wheel Prototype Classification Objective}

The first auxiliary objective imposes an absolute class reference. For a
sample \(i\), modality \(m\), and class \(c\), we define the prototype score as
the negative Poincar\'e distance scaled by temperature \(\tau_w>0\):
\begin{equation}
    s_{i,c}^{m}
    =-\frac{d_{\mathbb{D}}(z_i^{m},q_c)}{\tau_w}.
    \label{eq:prototype_score}
\end{equation}
The resulting prototype probability is
\begin{equation}
    \rho_{i,c}^{m}
    =\frac{\exp(s_{i,c}^{m})}
    {\sum_{k=1}^{C}\exp(s_{i,k}^{m})}.
    \label{eq:prototype_probability}
\end{equation}
We then minimize the mean cross-entropy over all modalities and valid
utterances:
\begin{equation}
    \mathcal{L}_{\mathrm{proto}}
    =-\frac{1}{|\mathcal{M}|\,|\mathcal{I}|}
    \sum_{m\in\mathcal{M}}
    \sum_{i\in\mathcal{I}}
    \log\rho_{i,y_i}^{m}.
    \label{eq:prototype_ce}
\end{equation}
This objective attracts each projected representation to the prototype of its
ground-truth emotion and simultaneously separates it from competing class
prototypes through the softmax normalization. Since the same prototype bank is
used for all modalities, \(\mathcal{L}_{\mathrm{proto}}\) also serves as a
weak cross-modal semantic alignment constraint without requiring paired
feature matching.

\subsubsection{Wheel-Structured CPCC Objective}

The prototype objective in \eqref{eq:prototype_ce} is pointwise: it encourages
correct classification but does not explicitly require the global arrangement
of utterances to preserve the relative distances between emotions. We address
this limitation with a cophenetic correlation coefficient (CPCC) objective.

The normalized shortest circular distance between classes \(a\) and \(b\) is
defined as
\begin{equation}
    \delta_{W}(a,b)
    =\frac{\arccos\!\left[
    \cos(\theta_a-\theta_b)\right]}{\pi}
    \in[0,1].
    \label{eq:wheel_class_distance}
\end{equation}
This target distinguishes adjacent, intermediate, and opposite emotions on the
Wheel instead of reducing every pair to the binary relation of same versus
different class.

Let
\(\mathcal{P}=\{(i,j):i,j\in\mathcal{I},\ i<j\}\)
be the set of distinct valid utterance pairs in the current mini-batch. For
each modality, we form a vector of geometric distances and a shared vector of
Wheel distances:
\begin{equation}
    g_{ij}^{m}=d_{\mathbb{D}}(z_i^{m},z_j^{m}),
    \qquad
    w_{ij}=\delta_W(y_i,y_j),
    \qquad(i,j)\in\mathcal{P}.
    \label{eq:cpcc_pair_distances}
\end{equation}
Denoting the corresponding pairwise means by
\(\bar{g}^{m}\) and \(\bar{w}\), let
\(\widetilde{g}_{ij}^{m}=g_{ij}^{m}-\bar{g}^{m}\) and
\(\widetilde{w}_{ij}=w_{ij}-\bar{w}\). Their Pearson correlation is
\begin{equation}
    \operatorname{corr}(g^{m},w)
    =
    \frac{\sum_{(i,j)\in\mathcal{P}}
    \widetilde{g}_{ij}^{m}\widetilde{w}_{ij}}
    {\max\!\left(
    \lVert\widetilde{g}^{m}\rVert_2
    \lVert\widetilde{w}\rVert_2,\epsilon\right)}.
    \label{eq:cpcc_correlation}
\end{equation}
The feature-level Wheel CPCC loss is
\begin{equation}
    \mathcal{L}_{\mathrm{CPCC}}
    =\frac{1}{|\mathcal{M}|}
    \sum_{m\in\mathcal{M}}
    \left[1-\operatorname{corr}(g^{m},w)\right].
    \label{eq:wheel_cpcc}
\end{equation}
When a mini-batch does not contain sufficient pairwise variation, the
corresponding term is set to zero for numerical stability. Minimizing
\eqref{eq:wheel_cpcc} encourages hyperbolic distances between utterances to
increase monotonically with their semantic separation on the Emotion Wheel.
Unlike a prototype-structure loss, this objective is applied directly to the
projected sample representations; the fixed prototypes themselves require no
additional structural regularization.

\subsubsection{Joint Optimization and Inference}

The complete training objective combines the original SDT loss with the two
geometry-aware regularizers:
\begin{equation}
    \boxed{
    \mathcal{L}
    =\mathcal{L}_{\mathrm{SDT}}
    +\lambda_p\mathcal{L}_{\mathrm{proto}}
    +\lambda_c\mathcal{L}_{\mathrm{CPCC}}}
    .
    \label{eq:complete_objective}
\end{equation}
The reference configuration uses
\(D=16\), \(r=0.75\), \(\tau_w=1\),
\(\lambda_p=0.1\), \(\lambda_c=0.05\), and
\(\gamma_1=\gamma_2=\gamma_3=1\). Both Wheel losses are active from the
first training epoch. Gradients from
\(\mathcal{L}_{\mathrm{proto}}\) and
\(\mathcal{L}_{\mathrm{CPCC}}\) update the modality projectors and propagate
into the SDT encoders; the fixed prototypes \(q_c\) receive no gradients.

The method is operated in an observation-only mode. In particular, the
historical anchor-bank typicality, consistency score, generic hyperbolic CPCC,
and prototype-logit fusion mechanisms do not modify either the SDT
distillation or the multimodal fusion weights. The prediction rule at test time
therefore remains
\begin{equation}
    \hat{y}_i
    =\underset{c\in\mathcal{Y}}{\arg\max}\;p_i[c],
    \label{eq:method_inference}
\end{equation}
where \(p_i\) is the original fused SDT distribution from
\eqref{eq:sdt_output}. This design makes the proposed branch an exclusively
training-time inductive bias: the prototype classifier can improve the
organization of modality representations without directly overriding the
stronger contextual multimodal classifier.

\section{Experimental Setup}
\label{sec:experimental_setup}

\section{Experimental Results}
\label{sec:experimental_results}

\section{Conclusion}
\label{sec:conclusion}

\end{document}
